Fix \(c\), \(k\), \(\ell\), and \(z\in\mathcal S_k\).  The coordinate sum of \(z\) is \(m_k\).  By \(\texttt{ass:bernoulli-label-iid}\), \(\texttt{ass:bernoulli-label-schedule-independence}\), and \(\texttt{ass:bernoulli-units}\),
\[
 \Pr(L_c=\ell,Z_c=z)
 =p_\ell\pi_\ell^{m_k}(1-\pi_\ell)^{n-m_k}
 =\frac{p_\ell B_{k\ell}}{M_k}. \tag{1}
\]
Summing (1) over \(\ell\) and using \(q=Bp\) gives
\[
 \Pr(Z_c=z)=\frac1{M_k}\sum_{\ell=1}^Kp_\ell B_{k\ell}=\frac{q_k}{M_k}. \tag{2}
\]
Every \(B_{k\ell}\) is positive because \(1\leq m_k\leq n-1\) and \(0<\pi_\ell<1\).  Since \(p\in\Delta_K\), at least one \(p_\ell\) is positive, whence \(q_k>0\).  Division by (2) is therefore valid, and Bayes' identity yields
\[
 \Pr(L_c=\ell\mid Z_c=z)=\frac{p_\ell B_{k\ell}}{q_k}. \tag{3}
\]
The right side of (2) is constant over the \(M_k\) elements of \(\mathcal S_k\).  Hence \(\Pr(Z_c\in\mathcal S_k)=q_k\) and \(Z_c\) is conditionally uniform on that slice.  Distinct \(m_k\)'s make the slices pairwise disjoint.  Finally the cluster labels, schedules, and conditional randomizations are independent across \(c\); thus the indicators of the \(K\) disjoint hit events and their complement are iid categorical with probabilities \((q_1,\ldots,q_K,1-\sum_kq_k)\).  Their totals have the asserted multinomial law.